\section{Алгоритм ро Полларда}

Алгоритм ро Полларда является общим вероятностным методом решения задачи дискретного логарифмирования в конечной циклической группе. Для задачи на эллиптических кривых его значение особенно велико, поскольку при корректном выборе параметров именно этот метод задаёт базовую оценку сложности атаки. Метод не использует разложение порядка группы и не опирается на специальные свойства кривой. Единственным существенным предположением является то, что задача рассматривается в подгруппе большого порядка, обычно простого (в последующих ).

Пусть точка $P \in E(\mathbb{F}_p)$ имеет порядок $n$, а точка $Q$ принадлежит подгруппе $\langle P\rangle$. Требуется найти число $k$, для которого $Q=kP$. Основная идея метода состоит в поиске двух различных представлений одной и той же точки группы в виде линейной комбинации точек $P$ и $Q$. Если удаётся найти коэффициенты $a',b',a'',b''$, для которых выполнено равенство
\[
a'P+b'Q=a''P+b''Q,
\]
то после подстановки $Q=kP$ получается соотношение
\[
(a'-a'')P=(b''-b')Q=(b''-b')kP.
\]
Так как точка $P$ имеет порядок $n$, отсюда следует линейное сравнение
\[
a'-a''\equiv (b''-b')k \pmod n.
\]
Если разность $b''-b'$ обратима по модулю $n$, то дискретный логарифм восстанавливается по формуле
\[
k\equiv (a'-a'')(b''-b')^{-1} \pmod n.
\]
Следовательно, вся задача сводится к эффективному поиску коллизии между двумя различными линейными представлениями одной и той же точки.

Наивный способ поиска коллизии состоял бы в генерации большого числа случайных точек вида $aP+bQ$ с последующим сохранением их в памяти до появления совпадения. По эффекту дней рождения коллизия в группе порядка $n$ ожидается примерно после $\sqrt{n}$ шагов. Однако такой подход требует хранения большого числа промежуточных значений. Алгоритм ро Полларда позволяет достичь той же асимптотической трудоёмкости при пренебрежимо малых требованиях к памяти.

Для этого строится псевдослучайная последовательность точек
\[
X_0,X_1,X_2,\dots,
\]
где каждая точка представляется в форме
\[
X_i=a_iP+b_iQ.
\]
Переход от $X_i$ к $X_{i+1}$ задаётся итерационной функцией, имитирующей случайное блуждание по подгруппе $\langle P\rangle$. Обычно группа разбивается на несколько подмножеств с помощью простой функции разбиения, зависящей, например, от нескольких битов координаты точки. Для каждого подмножества заранее выбирается некоторое приращение $R_j=\alpha_jP+\beta_jQ$. Если текущая точка $X_i$ попадает в подмножество с номером $j$, то следующий элемент последовательности вычисляется как
\[
X_{i+1}=X_i+R_j.
\]
Одновременно обновляются коэффициенты $a_i$ и $b_i$, чтобы в любой момент сохранялось равенство $X_i=a_iP+b_iQ$.

Так как подгруппа конечна, последовательность неизбежно начинает повторяться. После некоторого начального участка возникает цикл, и траектория принимает характерную форму, напоминающую греческую букву $\rho$, откуда и происходит название алгоритма. Если итерационная функция ведёт себя достаточно похоже на случайное отображение, то первая коллизия возникает примерно после $\sqrt{\pi n/2}$ шагов. Это и даёт оценку трудоёмкости метода порядка $O(\sqrt{n})$ групповых операций.

Для обнаружения коллизии при минимальном объёме памяти обычно используется алгоритм Флойда поиска цикла. Поддерживаются две последовательности, движущиеся по одной и той же траектории с разной скоростью: одна делает по одному шагу, а другая --- по два шага за итерацию. Как только их значения совпадают, фиксируется коллизия, после чего по сохранённым коэффициентам вычисляется искомый дискретный логарифм. Вероятность неудачи возникает только в случае, когда разность коэффициентов при $Q$ оказывается необратимой по модулю $n$; при простом $n$ эта ситуация имеет пренебрежимо малую вероятность.

Ещё одна существенная особенность метода состоит в возможности эффективного распараллеливания. Несколько вычислительных процессов могут независимо строить разные траектории случайного блуждания, используя одну и ту же итерационную функцию. Для обнаружения пересечений между траекториями удобно применять метод выделенных точек: процессор передаёт на центральный узел только точки, обладающие легко проверяемым отличительным свойством. Когда одна и та же выделенная точка обнаруживается повторно, из соответствующих коэффициентов восстанавливается дискретный логарифм. Такой подход обеспечивает практически линейное ускорение по числу процессоров и поэтому важен при оценке масштабируемых атак.

Для некоторых классов кривых метод может быть дополнительно ускорен за счёт использования легко вычислимых автоморфизмов группы точек. В этом случае несколько различных точек можно рассматривать как эквивалентные с точки зрения поиска коллизии, что уменьшает фактическое пространство поиска.